\documentclass[a4paper,12pt]{article}
\usepackage[T2A]{fontenc}
\usepackage[utf8]{inputenc}
\usepackage[russian]{babel}
\usepackage{amsmath,amssymb,mathtools}
\usepackage{helvet}
\renewcommand{\familydefault}{\sfdefault}
\usepackage{icomma}
\usepackage{geometry}
\geometry{left=20mm,right=20mm,top=20mm,bottom=21mm}
\usepackage{microtype}
\usepackage{tikz}
\usetikzlibrary{arrows.meta,positioning,shapes.geometric,calc}
\usepackage{tabularx,booktabs,array}
\usepackage{enumitem}
\usepackage{parskip}
\usepackage{fancyhdr}
\usepackage{setspace}
\usepackage{xcolor}
\usepackage[hidelinks]{hyperref}
\definecolor{accent}{HTML}{008080}
\definecolor{darkaccent}{HTML}{004D4D}
\definecolor{textgray}{HTML}{333333}
\definecolor{lightaccent}{HTML}{EDF5F3}
\setlength{\parindent}{0pt}
\setlength{\parskip}{5pt}
\onehalfspacing
\setlist[itemize]{leftmargin=7mm,itemsep=1pt,topsep=3pt}
\setlist[enumerate]{leftmargin=8mm,itemsep=9pt,topsep=5pt}
\renewcommand{\arraystretch}{1.28}
\pagestyle{fancy}
\fancyhf{}
\fancyhead[L]{\footnotesize\color{darkaccent}Финансовые задачи ЕГЭ}
\fancyhead[R]{\footnotesize\color{textgray}22.10.2026}
\fancyfoot[C]{\footnotesize\thepage}
\renewcommand{\headrulewidth}{0.3pt}
\setlength{\headheight}{15pt}
\newcommand{\mcheck}{\textcolor{accent}{\(\square\)}\hspace{3pt}}
\newcommand{\unit}[1]{\,\text{#1}}
\newcommand{\minorhead}[1]{\vspace{4pt}{\large\bfseries\color{darkaccent}#1}\par\vspace{0pt}}
\newcommand{\lessonhead}[1]{\section*{\color{darkaccent}#1}}
\begin{document}
\hypersetup{pageanchor=false}
\begin{titlepage}
\thispagestyle{empty}
\begin{center}
\vspace*{30mm}
{\color{accent}\Large МАТЕМАТИКА \quad\textbullet\quad ЕГЭ}\par
\vspace{18mm}
{\Huge\bfseries Чек-лист}\par
\vspace{6mm}
{\LARGE\color{darkaccent}\bfseries Финансовые задачи:}\par
\vspace{2mm}
{\LARGE\color{darkaccent}кредиты и вклады}\par
\vspace{17mm}
{\large Аннуитетная, дифференцированная,}\par
{\large произвольная и смешанная схемы}\par
\vspace{20mm}
{\large Матвей Горбачёв}\par
\vspace{3mm}
{\normalsize 22 октября 2026 года}\par
\vfill
{\small\color{textgray}Лёвин Артём Александрович}\par
{\small\color{textgray}эксклюзивно для Матвея Горбачёва}\par
\vspace{10mm}
\begin{tikzpicture}
\draw[accent!75,line width=0.9pt] (0,0) .. controls (2.3,0.6) and (5.4,-0.7) .. (8.0,0.15) .. controls (10.0,0.8) and (12.1,-0.4) .. (14,0.12);
\end{tikzpicture}
\end{center}
\end{titlepage}
\hypersetup{pageanchor=true}

\lessonhead{1. Как распознать финансовую схему}
\mcheck\textbf{Аннуитетная:} долг погашают \emph{равными платежами}. Каждый платёж включает проценты и часть основного долга; сама сумма платежа постоянна.

\mcheck\textbf{Дифференцированная:} основной долг уменьшается \emph{равными долями} (или каждый год на одну и ту же сумму). Проценты считают на уменьшающийся остаток, поэтому платежи обычно снижаются.

\mcheck\textbf{Произвольная:} условие задаёт остатки долга, отдельные выплаты или ограничения на них. Составляем таблицу движения долга.

\mcheck\textbf{Смешанная:} на разных этапах действуют разные правила. Например, первые три года платят только проценты, затем долг закрывают двумя равными платежами.

\minorhead{Обозначения и универсальный переход}
Пусть \(S\)~--- начальная сумма кредита, \(r\)~--- процентная ставка за период, \(k=1+\frac{r}{100}\)~--- повышающий коэффициент, \(B_i\)~--- остаток после \(i\)-го платежа, \(P_i\)~--- платёж.
\[
B_0=S,\qquad D_i=kB_{i-1},\qquad
\boxed{P_i=D_i-B_i},\qquad \boxed{B_i=kB_{i-1}-P_i}.
\]
Здесь \(D_i\)~--- долг \emph{после начисления процентов, до внесения платежа}. После последнего платежа обязательно \(B_n=0\).

\lessonhead{2. Аннуитет: равные платежи}
\mcheck\textbf{Формула для \(n\) равных платежей \(x\):}
\[
S k^n=x(k^{n-1}+k^{n-2}+\ldots+k+1),\qquad
\boxed{x=\frac{Sk^n}{1+k+\ldots+k^{n-1}}}.
\]
Она получается последовательным применением равенства \(B_i=kB_{i-1}-x\) и условия \(B_n=0\).

\textbf{Пример.} Кредит \(910\) тыс. руб. под \(20\%\) на три года погашают тремя равными платежами:
\[
x=\frac{910\cdot1{,}2^3}{1+1{,}2+1{,}2^2}
=\frac{1572{,}48}{3{,}64}=432\ \text{тыс. руб.}
\]
Общая сумма выплат: \(3\cdot432=1296\) тыс. руб., или \(1\,296\,000\) руб.

\minorhead{Вклад с пополнениями}
\mcheck\textbf{Деньги, внесённые раньше, приносят проценты дольше.} Вклад \(S\), открытый на три года под \(10\%\), и два пополнения \(a\) после первого и второго года дадут
\[
\boxed{F=Sk^3+ak^2+ak},\qquad k=1{,}1.
\]
\textbf{Пример.} Первоначально внесли \(100S\) тыс. руб., затем дважды по \(400\) тыс. руб.; итог \(1722{,}6\) тыс. руб. Тогда
\[
100S\cdot1{,}1^3+400\cdot1{,}1^2+400\cdot1{,}1=1722{,}6,
\qquad S=6.
\]
Первоначальный вклад равен \(600\) тыс. руб.

\lessonhead{3. Дифференцированная схема}
\mcheck\textbf{Каждый период погашают одинаковую часть основного долга:} \(S/n\). На начало периода \(i\) остаётся \(S-\frac{(i-1)S}{n}\). Поэтому
\[
\boxed{P_i=\frac{S}{n}+\frac{r}{100}\left(S-\frac{(i-1)S}{n}\right)}.
\]
\textbf{Пример.} Кредит \(2{,}5\) млн руб. на четыре года под \(15\%\). Ежегодно погашают по \(0{,}625\) млн руб. основного долга, а проценты равны \(0{,}375;\ 0{,}28125;\ 0{,}1875;\ 0{,}09375\) млн руб. Итого уплачено
\[
2{,}5+0{,}375+0{,}28125+0{,}1875+0{,}09375=3{,}4375\ \text{млн руб.}
\]

\lessonhead{4. Произвольная схема: таблица из пяти величин}
\mcheck\textbf{Универсальная последовательность для каждой строки:}
\[
\text{Начальный долг}\ \longrightarrow\ \text{Проценты}
\ \longrightarrow\ \text{Долг с процентами}
\ \longrightarrow\ \text{Платёж}\ \longrightarrow\ \text{Остаток}.
\]
\[
\boxed{I_i=\frac{r}{100}B_{i-1}},\quad
\boxed{D_i=B_{i-1}+I_i},\quad
\boxed{P_i=D_i-B_i}.
\]
\textbf{Пример.} Взяли \(750\) тыс. руб. на три года под \(12\%\). После первого и второго года остаток уменьшается на \(20\%\) относительно предыдущего года; после третьего~--- ноль.

\begin{center}
\small
\begin{tabular}{@{}c r r r r r@{}}
\toprule
Год & Начало & Проценты & С процентами & Платёж & Остаток\\
\midrule
1 & 750 & 90 & 840 & 240 & 600\\
2 & 600 & 72 & 672 & 192 & 480\\
3 & 480 & 57,6 & 537,6 & 537,6 & 0\\
\bottomrule
\end{tabular}
\end{center}
\vspace{-2pt}
\textit{Все суммы в таблице~--- в тысячах рублей.}

\[
\boxed{P_{\text{общ}}=240+192+537{,}6=969{,}6\ \text{тыс. руб.}=969\,600\ \text{руб.}}
\]
\textbf{Важно:} в последнем году погашаем \emph{весь} оставшийся долг с процентами. Уменьшение на \(20\%\) относилось лишь к первым двум годам.

\minorhead{Остатки можно задавать буквами}
Если исходная сумма неизвестна и обозначена \(S\), ставка \(5\%\) за период, а плановые остатки последовательно \(0{,}9S,\ 0{,}8S,\ 0{,}7S\), то первые строки таблицы:
\begin{center}
\small
\begin{tabular}{@{}c c c c c c@{}}
\toprule
Период & Начало & Проценты & С процентами & Платёж & Остаток\\
\midrule
1 & \(S\) & \(0{,}05S\) & \(1{,}05S\) & \(0{,}15S\) & \(0{,}9S\)\\
2 & \(0{,}9S\) & \(0{,}045S\) & \(0{,}945S\) & \(0{,}145S\) & \(0{,}8S\)\\
3 & \(0{,}8S\) & \(0{,}04S\) & \(0{,}84S\) & \(0{,}14S\) & \(0{,}7S\)\\
\bottomrule
\end{tabular}
\end{center}
\mcheck\textbf{Процент переплаты:}
\[
\boxed{q=\frac{P_{\text{общ}}-S}{S}\cdot100\%}.
\]
Фраза «на сколько процентов общая сумма выплат больше суммы кредита» требует именно этой формулы. Если процент начисляют на остаток и нет дополнительных сборов, удобно проверить себя: общая переплата равна сумме начисленных процентов.

\lessonhead{5. Ограничения на платежи и целый ответ}
\mcheck\textbf{Различайте} условия «\emph{каждый} платёж меньше \(L\)» и «\emph{сумма всех} платежей меньше \(L\)».

\textbf{Пример.} Ставка \(20\%\), остатки \(S\to0{,}8S\to0{,}4S\to0\). Платежи:
\[
P_1=0{,}4S,\qquad P_2=0{,}56S,\qquad P_3=0{,}48S.
\]
Если \emph{каждый} платёж меньше \(5\) млн руб., то одновременно
\[
\begin{cases}
0{,}4S<5,\\
0{,}56S<5,\\
0{,}48S<5,
\end{cases}
\quad\Longleftrightarrow\quad
\begin{cases}
S<12{,}5,\\
S<\dfrac{125}{14}\approx8{,}929,\\
S<\dfrac{125}{12}\approx10{,}417.
\end{cases}
\]
При \(S\in\mathbb N\) наибольшее допустимое \(S=8\) млн руб.

\mcheck\textbf{Если ограничена общая сумма,} сперва сложите платежи: \(P_{\text{общ}}=(0{,}4+0{,}56+0{,}48)S=1{,}44S\), затем решайте \(1{,}44S<L\).

\lessonhead{6. Смешанная схема: сначала проценты, потом аннуитет}
\mcheck\textbf{Разделяйте кредит на этапы.} Пока выплачиваются только проценты, основной долг остаётся неизменным. Равные платежи в конце удобно обозначить одной буквой \(x\).

\textbf{Пример.} Кредит \(3{,}3\) млн руб. под \(20\%\) на пять лет. В первые три года платят лишь проценты, в последние два~--- по \(x\) млн руб. Первый платёж \(0{,}2\cdot3{,}3=0{,}66\) млн руб. В четвёртый год перед платежом должны \(3{,}96\) млн руб., после него~--- \(3{,}96-x\). В пятый год кредит полностью закрывают:
\[
1{,}2(3{,}96-x)-x=0,
\qquad x=2{,}16\ \text{млн руб.}
\]
Разница последнего и первого платежей: \(2{,}16-0{,}66=1{,}5\) млн руб.

\minorhead{Контроль перед записью ответа}
\begin{itemize}
\item Проверить, с какой даты начисляют проценты и в какой момент вносят платёж.
\item В последнем периоде поставить остаток \(0\); проверить, что долг действительно погашен.
\item Проверить, требуется ли \emph{платёж}, \emph{общая сумма}, \emph{переплата} или \emph{разница платежей}.
\item Перевести тысячи или миллионы рублей в рубли, если это указано в вопросе.
\item При строгом знаке «\(<\)» граничное значение не включать; условие \(S\in\mathbb N\) учитывать отдельно.
\end{itemize}

\clearpage
\thispagestyle{plain}
\begin{center}
{\Large\bfseries\color{darkaccent}Алгоритм: определяем схему кредитования}\par
\vspace{6mm}
{\normalsize Читайте ключевую формулировку условия, затем выбирайте модель расчёта.}
\end{center}
\vspace{7mm}
\begin{center}
\begin{tikzpicture}[
 node distance=10mm and 8mm,
 base/.style={draw=darkaccent,line width=.9pt,align=center,inner sep=5pt,font=\small,text width=55mm},
 terminus/.style={base,ellipse,text width=46mm},
 question/.style={base,diamond,aspect=1.8,text width=35mm,inner xsep=3pt,inner ysep=6pt},
 action/.style={base,rounded corners=3pt},
 arr/.style={-Latex,line width=1pt,color=accent}
]
\node[terminus] (start) {Прочитайте способ погашения долга};
\node[question,below=of start] (a) {Все платежи равны?};
\node[action,right=8mm of a] (an) {Аннуитет: один платёж \(x\); формула равных платежей};
\node[question,below=15mm of a] (d) {Основной долг уменьшается равными долями?};
\node[action,right=8mm of d] (df) {Дифференцированная: одинаковая доля тела кредита};
\node[action,below=12mm of d] (arb) {Произвольная или смешанная: таблица долга, \(D_i=kB_{i-1}\), \(P_i=D_i-B_i\); в конце \(B_n=0\)};
\draw[arr] (start)--(a);
\draw[arr] (a.east)--node[above,font=\small]{да}(an.west);
\draw[arr] (a.south)--node[left,font=\small]{нет}(d.north);
\draw[arr] (d.east)--node[above,font=\small]{да}(df.west);
\draw[arr] (d.south)--node[left,font=\small]{нет}(arb.north);

\end{tikzpicture}
\end{center}

\clearpage
\lessonhead{Тренировочный блок \texorpdfstring{\(\cdot\)}{} Путь А}
\textit{10 заданий. Записывайте основную формулу или таблицу; в условии с целым \(S\) обязательно проверяйте строгость неравенства.}
\begingroup
\setstretch{1.20}
\begin{enumerate}[label=\textbf{\arabic*.},itemsep=5pt]
\item Вклад открыли на три года под \(10\%\) годовых, внеся \(100S\) тыс. руб., где \(S\)~--- натуральное число. После первого и второго года вклад пополняли на \(200\) тыс. руб. Итоговая сумма составила \(1260{,}6\) тыс. руб. Найдите \(S\).
\item Кредит \(210\) тыс. руб. взят на два года под \(10\%\) годовых. Долг погашен двумя равными ежегодными платежами. Найдите \emph{общую сумму} платежей в рублях.
\item Кредит \(800\) тыс. руб. под \(10\%\) на четыре года возвращают равными долями основного долга. Проценты ежегодно начисляют на оставшуюся задолженность. Найдите общую сумму выплат в рублях.
\item Кредит \(1\) млн руб. взят на два года под \(10\%\). После первого года остаток должен быть на \(30\%\) меньше первоначального долга; после второго кредит закрывают. Найдите сумму двух платежей в рублях.
\item Кредит \(900\) тыс. руб. выдан на три года под \(10\%\). После первого и второго года остаток долга должен составлять \(80\%\) остатка перед соответствующим начислением процентов; в конце третьего года долг погашается полностью. Найдите общую сумму выплат в рублях.
\item Кредит \(S\) выдан на три года под \(10\%\). После первого и второго года остаток составляет \(80\%\) долга на начало соответствующего года; после третьего~--- ноль. На сколько процентов все платежи вместе превышают сумму \(S\)?
\item Кредит \(1{,}5\) млн руб. под \(5\%\) за месяц погашают за четыре месяца. Остатки после первых трёх платежей составляют соответственно \(0{,}9S\), \(0{,}8S\) и \(0{,}7S\), где \(S=1{,}5\) млн руб. После четвёртого платежа долг равен нулю. Найдите общую сумму платежей в рублях.
\item Кредит \(S\) млн руб., где \(S\)~--- натуральное число, взят на три года под \(20\%\). Остатки после платежей равны \(0{,}8S\), \(0{,}4S\), \(0\). Найдите \emph{наибольшее} \(S\), при котором \emph{каждый} из трёх платежей строго меньше \(6\) млн руб.
\item При тех же ставке и схеме остатков, что в задании 8, найдите наибольшее натуральное \(S\), для которого \emph{общая сумма всех трёх платежей} строго меньше \(13\) млн руб.
\item Кредит \(2{,}75\) млн руб. взят на пять лет под \(20\%\) годовых. В первые три года платят только начисленные проценты, а в последние два года~--- равными платежами, полностью закрывая долг. На сколько рублей последний платёж больше первого?
\end{enumerate}
\endgroup

\clearpage
\lessonhead{Ответы к тренировочному блоку}
\begin{center}
\renewcommand{\arraystretch}{1.55}
\begin{tabularx}{\linewidth}{@{}c X@{}}
\toprule
\textbf{№} & \textbf{Ответ}\\
\midrule
1 & \(S=6\).\\
2 & \(242\,000\) руб.\\
3 & \(1\,000\,000\) руб.\\
4 & \(1\,170\,000\) руб.\\
5 & \(1\,119\,600\) руб.\\
6 & \(24{,}4\%\).\\
7 & \(1\,755\,000\) руб.\\
8 & \(S=10\) (млн руб.).\\
9 & \(S=9\) (млн руб.).\\
10 & \(1\,250\,000\) руб.\\
\bottomrule
\end{tabularx}
\end{center}
\vspace{7mm}
\textbf{Самопроверка.} В задачах 8 и 9 используются одни и те же платежи:
\[
P_1=0{,}4S,\qquad P_2=0{,}56S,\qquad P_3=0{,}48S,
\qquad P_{\text{общ}}=1{,}44S.
\]
Однако в задаче 8 ограничен \emph{каждый платёж}, а в задаче 9~--- \emph{их сумма}. Поэтому ответы могут различаться.
\vfill
\begin{center}
{\small\color{textgray}Лёвин Артём Александрович эксклюзивно для Матвея Горбачёва}
\end{center}
\end{document}
